\clearpage
\section{Runtime Instruction Examples}

This section shows how runtime support code can combine architectural discovery instructions with
state-management instructions. The examples are schematic C helpers, not ABI-mandated entry points.

\subsection{CPUID-Defined Program Context Save}
SAVE and RESTORE use a CPUID-described memory image. A runtime that wants to save program-visible register and
extension state first queries the required save-area size and component descriptors, then supplies a
4 KiB-aligned buffer to SAVE.

The continuation state that is not part of the SAVE area, such as a scheduler or setjmp return point, is managed
by the runtime layer that calls this helper.

The two target hooks below are ordinary runtime entry points implemented with the architectural \texttt{CPUID Rn}
and \texttt{SAVE \textless{}ea\textgreater{}} instructions.

\par\begingroup\footnotesize
\begin{verbatim}
#include <stddef.h>
#include <stdint.h>

#define BR_CPUID_CLASS_IMPL       0x00000002u
#define BR_CPUID_SAVE_AREA_LAYOUT 0x0004u
#define BR_SAVE_ALIGNMENT         4096u

struct br_save_area_info {
    uint64_t size;
    uint16_t header_size;
    uint16_t component_count;
    uint16_t bitmap_words;
};

struct br_save_component {
    uint16_t id;
    uint16_t flags;
    uint32_t offset;
};

struct br_program_context {
    void *save_area;
    uint64_t save_area_size;
    uint16_t header_size;
    uint16_t component_count;
    uint16_t bitmap_words;
};

uint64_t bedrock_cpuid(uint64_t selector);
void bedrock_save(void *save_area);

static uint64_t br_cpuid_selector(uint32_t cls,
                                  uint16_t leaf,
                                  uint16_t index)
{
    return ((uint64_t)cls << 32) |
           ((uint64_t)leaf << 16) |
           (uint64_t)index;
}

static struct br_save_area_info br_query_save_area(void)
{
    uint64_t r0 = bedrock_cpuid(
        br_cpuid_selector(BR_CPUID_CLASS_IMPL,
                          BR_CPUID_SAVE_AREA_LAYOUT, 0));
    uint64_t r1 = bedrock_cpuid(
        br_cpuid_selector(BR_CPUID_CLASS_IMPL,
                          BR_CPUID_SAVE_AREA_LAYOUT, 1));
    struct br_save_area_info info;

    info.size = r0;
    info.header_size = (uint16_t)(r1 & 0xffffu);
    info.component_count = (uint16_t)((r1 >> 16) & 0xffffu);
    info.bitmap_words = (uint16_t)((r1 >> 32) & 0xffffu);
    return info;
}

static void br_query_save_components(struct br_save_component *out,
                                     uint16_t count)
{
    for (uint16_t i = 0; i != count; ++i) {
        uint64_t r = bedrock_cpuid(
            br_cpuid_selector(BR_CPUID_CLASS_IMPL,
                              BR_CPUID_SAVE_AREA_LAYOUT,
                              (uint16_t)(2u + i)));
        out[i].id = (uint16_t)(r & 0xffffu);
        out[i].flags = (uint16_t)((r >> 16) & 0xffffu);
        out[i].offset = (uint32_t)(r >> 32);
    }
}

static uintptr_t br_align_up(uintptr_t value, uintptr_t alignment)
{
    return (value + alignment - 1u) & ~(alignment - 1u);
}

int br_save_program_context(struct br_program_context *ctx,
                            struct br_save_component *components,
                            uint16_t component_capacity,
                            void *storage,
                            size_t storage_size)
{
    struct br_save_area_info info = br_query_save_area();
    uintptr_t raw = (uintptr_t)storage;
    uintptr_t base = br_align_up(raw, BR_SAVE_ALIGNMENT);
    size_t leading = (size_t)(base - raw);

    if (leading > storage_size)
        return -1;
    if (info.size > (uint64_t)(storage_size - leading))
        return -1;
    if (info.component_count > component_capacity)
        return -1;

    br_query_save_components(components, info.component_count);

    ctx->save_area = (void *)base;
    ctx->save_area_size = info.size;
    ctx->header_size = info.header_size;
    ctx->component_count = info.component_count;
    ctx->bitmap_words = info.bitmap_words;

    bedrock_save(ctx->save_area);
    return 0;
}
\end{verbatim}
\endgroup\par

The component array is not part of the SAVE image. It is ordinary runtime metadata derived from CPUID so the
context owner can inspect, copy, or migrate extension components without hard-coding implementation-specific
offsets.
